
\section{AA-01: Selection rules and radiative lifetime}
\source{Week 9 lectures; Week 10 spectroscopy; 2026 Assignment 4,
Question 3 (sodium lifetimes).}
\subsection*{Step 1: decide whether the matrix element can exist}
For an E1 transition, the dipole operator \(d_q\)
is a rank-1 tensor of \emph{odd} parity:
\[
\langle J'm'|d_q|Jm\rangle=
(-1)^{J'-m'}
\begin{pmatrix}J'&1&J\\-m'&q&m\end{pmatrix}
\langle J'\Vert d\Vert J\rangle .
\]
Thus \(m'=m+q\), \(|J-J'|\le1\le J+J'\),
and parity must change. For one electron in the
LS-coupled limit, E1 orbital selection gives
\(\Delta l=\pm1\). This predicts, for example,
\[
3p\to3s\ \text{E1 allowed},\quad
4s\to3p\ \text{E1 allowed},\quad
3d\to3p\ \text{E1 allowed},
\]
but \(3d\to3s\) is \emph{not} E1 allowed because
\(\Delta l=-2\). It may be possible by higher multipoles.
Whether a decay is physically available also requires
the final state to lie below the initial energy.

\subsection*{Step 2: separate angular and radial factors}
For \(\psi_{nlm}=R_{nl}(r)Y_{lm}(\Omega)\),
the dipole amplitude separates into
\[
\langle f|r_q|i\rangle=
\underbrace{\int_0^\infty
R_f^*(r)R_i(r)r^3\,dr}_{\text{radial overlap}}\!
\underbrace{\int Y_{l_fm_f}^*(\Omega)
C_q^{(1)}(\Omega)Y_{l_im_i}(\Omega)\,d\Omega}_{
\text{angular / selection factor}} .
\]
The reduced radial convention \(u(r)=rR(r)\)
uses \(\int u_f^*r u_i\,dr\) instead.
Never confuse \(R\) and \(u\).

\subsection*{Step 3: convert an E1 dipole into a rate}
For one specified transition with correctly summed
polarisation components,
\[
\boxed{A_{if}=\frac{\omega_{if}^3}{3\pi\epsilon_0\hbar c^3}
|\langle f|\mathbf d|i\rangle|^2},\quad
\omega_{if}=\frac{E_i-E_f}{\hbar}.
\]
If the excited state has several allowed lower states,
\(\Gamma_i=\sum_f A_{if}\) and
\(\boxed{\tau_i=1/\Gamma_i}\).
\textbf{Illustrative calculation (not measured Na data):}
take a single 589 nm transition with
\(|\mathbf d|=2ea_0\). Using
\(a_0=5.29177\times10^{-11}\) m and
\(\omega=2\pi c/\lambda\) gives
\[
A\simeq3.97\times10^7\ {\rm s}^{-1},\qquad
\boxed{\tau\simeq25.2\ {\rm ns}}
\]
if no other decay channels contribute.
For actual sodium \(3p,4s,3d\), calculate the
state-dependent dipole integrals and sum all
E1-allowed downward channels as the assignment asks.

\subsection*{Step 4: contrast E1, M1 and E2}
E1: rank 1, odd parity; M1: rank 1, even parity;
E2: rank 2, even parity. Rank-2 angular momentum
demands \(|J-J'|\le2\le J+J'\). In particular,
an \(s\to d\) orbital transition can be E2,
whereas \(s\to p\) cannot be E2 by parity.
\checkmethod{A nonzero 3-j symbol proves only angular
compatibility. The reduced matrix element may still vanish;
the full lifetime requires all permitted lower channels.}
